\documentclass[nonacm]{acmart}

\AtBeginDocument{%
  \providecommand\BibTeX{{%
    Bib\TeX}}}

\setcopyright{none}

\usepackage{array}
\usepackage{mathtools}
\usepackage{stmaryrd}
\usepackage{mathpartir}
\usepackage{listings}
\usepackage{xcolor}

\lstset{
  basicstyle=\ttfamily\small,
  breaklines=true,
  columns=fullflexible,
  keepspaces=true,
  frame=single,
  showstringspaces=false
}

\newcommand{\FV}{\mathsf{FV}}
\newcommand{\dom}{\mathsf{dom}}
\newcommand{\Nat}{\mathsf{Nat}}
\newcommand{\Bool}{\mathsf{Bool}}

\begin{document}

\title{Constraint Typing for Self-Application}

\author{Runbang Yan}
\email{rbyan25@stu.pku.edu.cn}
\affiliation{%
  \institution{School of Computer Science, Peking University}
  \city{Beijing}
  \country{China}
}

\renewcommand{\shortauthors}{Yan}

\maketitle

\section{Constraint Typing Derivation}

Consider the term
\[
  \lambda x:X.\ x\ x .
\]
In the body, both occurrences of $x$ have type $X$ under the context
$x:X$.  Applying rule \textsc{CT-App}, introduce a fresh type variable
$Y$ for the result of the application.  The application generates the
constraint that the type of the left occurrence of $x$ must be an arrow
from the type of the argument occurrence of $x$ to $Y$:
\[
  X = X \to Y .
\]
Therefore the constraint typing derivation is:
\[
\inferrule*[Right=\textsc{CT-Abs}]
  {
    \inferrule*[Right=\textsc{CT-App}]
      {
        \inferrule*[Right=\textsc{CT-Var}]
          {x:X \in x:X}
          {x:X \vdash x : X \mid_{\emptyset} \emptyset}
        \\
        \inferrule*[Right=\textsc{CT-Var}]
          {x:X \in x:X}
          {x:X \vdash x : X \mid_{\emptyset} \emptyset}
        \\
        Y \notin \{X\}
      }
      {x:X \vdash x\ x : Y \mid_{\{Y\}} \{X = X \to Y\}}
  }
  {
    \emptyset \vdash \lambda x:X.\ x\ x : X \to Y
      \mid_{\{Y\}} \{X = X \to Y\}
  } .
\]

\section{Unification with Finite Types}

The generated constraint set is
\[
  C = \{X = X \to Y\}.
\]
It is not unifiable in the ordinary finite type language.  Any unifier
$\sigma$ would have to satisfy
\[
  \sigma(X) = \sigma(X) \to \sigma(Y).
\]
Thus $\sigma(X)$ would need to be equal to a proper arrow type containing
itself on the left.  This cannot be represented by a finite type.
Equivalently, the standard unification algorithm rejects the equation by
the occurs check because
\[
  X \in \FV(X \to Y).
\]

\section{Unification with Recursive Types}

If the occurs check is replaced by the recursive-type case from the
slides, the equation
\[
  X = X \to Y
\]
matches the branch where $S=X$ and $X \in \FV(T)$, with
$T = X \to Y$.  Hence the algorithm computes
\[
\begin{aligned}
  \mathsf{unify}(\{X = X \to Y\})
    &=
      \mathsf{unify}([X \mapsto \mu X.\,X \to Y]\emptyset)
      \circ [X \mapsto \mu X.\,X \to Y] \\
    &=
      [X \mapsto \mu X.\,X \to Y].
\end{aligned}
\]
Under the usual unfolding equality for recursive types,
\[
  \mu X.\,X \to Y \equiv (\mu X.\,X \to Y) \to Y,
\]
so this substitution satisfies the original constraint.

Applying this substitution to the type assigned by the constraint
derivation gives
\[
  \sigma(X \to Y)
  =
  (\mu X.\,X \to Y) \to Y.
\]
Thus, with recursive types, the algorithm accepts the self-application
and returns the substitution
\[
  [X \mapsto \mu X.\,X \to Y],
\]
leaving $Y$ unconstrained.

\end{document}
